\begin{examplebox}[Taking Annotated Theorems as Hints: Example]
\textbf{Problem:} Let $a, b, c$ be the lengths of the sides of a triangle, such that $a + b + c = 3$. Determine the largest constant $C$ such that the following inequality holds for all $a, b, c$:
$$
a^{2} + b^{2} + c^{2} + \frac{4 a b c}{3} \geq C.
$$
\\
\textbf{Related Theorem:}\\ 
\correct{\textbf{Theorem Zdravko (Mean inequalities)}}: Let $a_{1}, a_{2}, \ldots, a_{n}$ be positive real numbers. The numbers
$$
\begin{aligned}
& Q M=\sqrt{\frac{a_{1}^{2}+a_{2}^{2}+\cdots+a_{n}^{2}}{n}}, \quad A M=\frac{a_{1}+a_{2}+\cdots+a_{n}}{n} \\
& G M=\sqrt[n]{a_{1} a_{2} \cdots a_{n}} \quad \text { and } \quad H M=\frac{n}{\frac{1}{a_{1}}+\frac{1}{a_{2}}+\cdots+\frac{1}{a_{n}}}
\end{aligned}
$$
are called the quadratic, arithmetic, geometric and harmonic mean for the numbers $a_{1}, a_{2}, \ldots, a_{n}$, respectively, and we have
$$
\correct{Q M \geq A M \geq G M \geq H M}
$$

Equalities occur if and only if $a_{1}=a_{2}=\cdots=a_{n}$.
\\\\
\textbf{Response by o4-mini without any Theorem Hint (Incorrect Response \xmark):}\\
Let $F(a, b, c) = a^2 + b^2 + c^2 + \frac{4}{3}abc$.\\
\wrong{\textbf{1. Equilateral case:} Let \( a = b = c = 1 \)}. Then \( a + b + c = 3 \), and
\[
F(1,1,1) = 1^2 + 1^2 + 1^2 + \frac{4}{3} \cdot 1 \cdot 1 \cdot 1 = 3 + \frac{4}{3} = \frac{13}{3}
\]

\wrong{\textbf{2. Degenerate case:} Let  \( a = b = 0.75,\ c = 1.5 \)}. Then,
\[
F(0.75, 0.75, 1.5) = 0.75^2 + 0.75^2 + 1.5^2 + \frac{4}{3} \cdot 0.75 \cdot 0.75 \cdot 1.5 = 4.5
\]

Thus, on the boundary we have \( F(a,b,c) \geq 4.5 \).\\

\wrong{A more detailed check of the two-equal-sides slice \( a = b = x,\ c = 3 - 2x \) shows that} \( F(x) \) attains its unique minimum \( 13/3 \) at \( x = 1 \). No other interior configuration yields a smaller value. Therefore, the answer is \boxed{C = \frac{13}{3}}.
\\

% \textbf{(Continued)}

% \end{examplebox}

% \begin{examplebox}[Taking Annotated Theorems as Hints: Example (Continued)]

\textbf{Response by o4-mini with a Theorem Hint (Correct Response \cmark):}
Set \(a+b+c=3\).  \correct{By the quadratic–arithmetic mean inequality (QM–AM) in \textbf{Mean inequalities theorem}},
\[
\frac{a^2+b^2+c^2}{3}\ge\Bigl(\frac{a+b+c}{3}\Bigr)^2=1
\quad\Longrightarrow\quad
a^2+b^2+c^2\ge3.
\]
\correct{By the AM–GM inequality in \textbf{Mean inequalities theorem}},
\[
\sqrt[3]{abc}\le\frac{a+b+c}{3}=1
\quad\Longrightarrow\quad
abc\le1.
\]
Hence
\[
a^2+b^2+c^2+\frac{4abc}{3}
\;\ge\;3+\frac{4\cdot1}{3}
\;=\;3+\frac43
\;=\;\frac{13}{3}.
\]
Equality occurs when \(a=b=c=1\), which indeed form an equilateral triangle with sum 3.  Thus the largest constant \(C\) for which the inequality holds for all triangles is
\[
C=\frac{13}{3}.
\]
The answer is \boxed{C=13/3}.
\\\\
\textbf{Expert Comment:} 
In the version without the theorem, the model relies on toy-case evaluations—specifically symmetric and degenerate configurations—to guess the correct bound \wronghl. In addition, its reasoning includes a clear logical gap: it claims that “a more detailed check of the two-equal-sides slice $a = b = x,\ c = 3 - 2x$” shows the minimum occurs at $x = 1$, but provides no explanation or computation to support this claim \wronghl. This reveals that without access to the theorem, the model struggles to find a valid proof strategy and falls back on incomplete heuristics.
\\\\
Once the relevant inequality theorems are provided, the model quickly applies the QM–AM and AM–GM inequalities from the \correct{\textbf{Mean inequalities theorem}} correctly \correcthl. It uses these tools to derive a general lower bound valid for all triangles, leading rigorously to the correct constant $C = \frac{13}{3}$. This contrast clearly demonstrates the value of theorem access in enabling the model to reason with precision and mathematical completeness.

\end{examplebox}